\chapter{反抗小圆}
名字拥有力量。

焰把哪些人拉进自己的迷宫，其实有迹可循。
魔法少女们最容易理解：
她们是焰生命中的核心人物，至少在最近的一条时间线里，也是她的同伴。
小圆的家人稍显绕远，但他们对小圆重要，而小圆对焰重要，因此仍说得过去。
仁美和恭介更绕一层，不过仁美之于小圆，恭介之于仁美，都构成了足够的关联。
真正耐人寻味的是，焰为什么要把和子老师和中泽也带进来：
一个是班主任，一个看上去只是普通同学。
但电影确实如此安排。
数学课上，焰开始怀疑身边人的真实性时，中泽和其他魔法少女一样，是仅有的仍保有正常面孔的人。
迷宫最终崩解时，中泽与和子老师\footnote{%
    译注：译者去考察时，只发现一人在椅子上闭目。
}也都昏倒在沙发上。

答案要到民俗学里找：
传说中，女巫若知道一个人的名字，便能借这个名字对其施法\cite{waite2006black-magic}。
再看焰带进迷宫的名单，并把整个动画一并纳入视野，就会发现这份名单同时符合另一个标准：
名单中人物的名字都曾被说出或出现在画面上，并且焰有理由知道这些人的姓名。
小圆父亲的名字固然没有在正片中被说出口，但小圆称他为“爸爸”，而全剧中又没有其他仍在世的父亲出场，所以“爸爸”已经足以承担名字的功能。

名字之所以有力量，是因为在魔法逻辑中，或者说在很大程度上也就是在叙事逻辑中，能指与所指并不分离。
名字在某种意义上就是被命名之物，因此，操纵名字就是操纵事物本身\cite{waite2006black-magic}。
于是，如果两样东西拥有同一个名字，它们便必然在某种意义上是同一的，其中一个可以替代另一个。

绕了这么一大圈，其实是为了说明：
焰抓住小圆的双臂，把“作为少女的小圆”从“作为抽象概念的小圆”中撕扯出来时，这不仅是焰对小圆的反叛，也是《叛逆》对《魔圆》本身的反叛。

电影反叛这部动画，又有何不可呢？
在很久以前，一个人想讲故事，就站起来开口讲。
口头讲述的故事转瞬即逝，每一次讲述都会有所变化。
传统当然存在，但叙事者可以相信，自己的创造性偏离并不会被听众视为错误或背叛，反而会被当成高明表演的即兴润色。

识字教育的普及给这种讲故事方式带来了致命一击，广播、电影和电视则完成了后续工作，让其轰然倒塌。
这种讲述并未彻底消失，毕竟没有哪种艺术形式会真正死亡，但它多半只剩下猎奇式的存在：
人们在文艺复兴节上驻足围观，或带孩子去公共图书馆听一场活动。
大多数时候，我们想要故事，拿起的是已经包装好的东西：
一本书，一张 DVD。

问题也由此出现：
当作者想讲一个故事系列，并反复使用同一批角色、同一套设定时，该怎么办？
作者想探索，想创造，这在口述故事时代，原本不成问题。
没有人会坚持《列那狐传奇》\footnote{%
    译注：《列那狐传奇》是中世纪集体口传的讽刺故事合集，由法国玛·阿希·季诺夫人整理成书，故各篇互有矛盾，并不要求设定严格统一。
}的每个故事都必须彼此吻合，也不会有人抱怨：
“宙斯诱惑勒达时明明是天鹅，怎么这回又变成金雨了\footnote{%
    译注：作者故意将两则希腊神话混为一谈——
    宙斯化作天鹅引诱勒达，化作金雨引诱达那厄，本非同一故事。
    作者以此论证口述传统中故事随讲述而变异，听众不苛求前后统一。
}？”
说到底，如果勒达的故事本就可以随着每一次讲述而改变，那么当你听到的是另一个故事时，又为什么要求它依旧保持同一副面貌？

口述故事是活物，会生长，会变化。
相反，写在纸上或拍在银幕上的故事已经死亡，被钉死在页面和影像之中，不能改变，也不能生长，只能永远停留在某一次讲述完成的形态。
观众可以改变，可以成长，因此他们对故事的理解也会随时间变动；
可故事的真正创作者却反而被剥夺了这种成长的权利。
即便在续作创作中，观众，尤其是以此闻名的“极客”与邪典\footnote{
    译注：邪典（Cult）是指风格独特、题材偏锋、具有争议性的小众艺术作品，通常不以市场为导向，但在特定群体中拥有狂热粉丝。
}受众，仍要求“连贯性”。
换言之，他们要求创作者效忠于已经死去的故事所建立的暴政。
更令人惊讶的是，竟没有多少创作者起而反叛。

于是，《叛逆》对动画的致敬可谓口惠而实不至。
动画中的事件在电影里显然都发生过，并获得了连贯性崇拜者所谓的“尊重”：
那条“不得与前作事件相矛盾”的律法，仅遵其字面，而未顾其精神。
电影甚至仿照了动画的结构，清晰地分为三段：
第一段伪装成一部“正常”的魔法少女作品，直到这层伪装在与麻美的激烈冲突中骤然崩塌。
第二段（当然，电影中这一段与第一段的重叠程度比动画更高）则跟随一名魔法少女，让她逐渐意识到自己一直在对抗的正是自己，而她从一开始就是魔女。
最后，第三段展开一场对抗城市规模魔女的大战，战后现实被改写，新的秩序随之建立。
% 英文原文：a tremendous battle against a city-scale witch

然而，动画跟随的是小圆，电影跟随的是焰，差异正在此处。
小圆耐心、谨慎，却又极其乐观。
她直到动画终局才行动，可一旦行动，便异常果决，意在彻底终结她所认定的宇宙根本问题。
焰则愤世嫉俗、执拗、充满对抗性。
她一次又一次把自己抛入冲突，直到最后，她以镜像般的方式回应小圆的行动：
她创造出一个世界，在那里，焰最大的难题，也就是小圆的自我牺牲冲动，必须被永无止境地处理下去。

这并不意味着在电影的大部分时间里，焰都自觉地反抗着小圆建立的秩序。
起初的焰确实仍像动画中那样，处在质疑并扰乱现状的位置上。
但这一次，仍由麻美代表并守护的现状，却是焰自己的梦境。
焰想恢复的是她记忆中的世界，也就是动画中的世界。
她真正开始有意识地反叛，是在小圆告诉她，与所爱之人分离会令她无比痛苦之后。
换句话说，是在焰意识到小圆的自我牺牲确实包含着真实的牺牲之后。
对焰而言，小圆的牺牲当然不可想象，也不可原谅，哪怕作出这一选择的正是小圆本人。

即便如此，焰也直到电影末段才真正有机会去否定小圆的牺牲。
在此之前，电影中真正反抗小圆，并因此反抗《魔圆》的角色，是丘比。
它策划了整个局面，试图夺取圆环之理的控制权，让魔女重新回到世界。
这里值得注意的是，丘比在许多方面都是一个极不讨喜的作者替身。
因此，它对小圆的反叛，也就可以理解为创作者借《叛逆》对《魔圆》发起的反叛。

不过，丘比的反叛并不意外，它毕竟是动画中的反派。
一个毫无悔意且仍在续作中登场的反派，至少会试图重拾自己的反派职能。
焰则不同。
她对小圆的忠诚强烈到近乎执念，因此电影必须格外细致地铺陈她反叛的每一个条件。
她有动机：
花田中与小圆的对话，使她意识到自己“绝不该允许”小圆牺牲自己。
她有启示：
丘比揭示小圆仍然可以选择重新进入世界，而沙耶香则揭示，小圆的佛性，即她作为圆环之理的记忆与力量，可以由他者暂时保存。
她也有机会：
当小圆降临，准备取走她的生命，阻止她在“现实世界”成为魔女时，正如丘比所说，凡可被感知者，皆可被干涉。

于是，焰以恶魔的形象起身，把“神”从天堂撕下，拖回尘世。
她是终极的坏女孩：引用《失乐园》\footnote{%
    译注：《失乐园》是英国政治家、学者约翰·弥尔顿创作的史诗，讲述了叛逆之神撒旦因为反抗上帝的权威被打入地狱，却毫不屈服，为复仇寻至伊甸园。
    亚当与夏娃受被撒旦附身的蛇引诱，偷吃了上帝明令禁食的知识树果实。
    最终，撒旦及其同伙遭谴，全都变成蛇，亚当与夏娃则被逐出伊甸园。
}的涂鸦，以及那些吟诵尼采、投掷番茄的使魔，都把她塑造成撒旦。
她主动认领了“恶魔”、“邪恶”、“敌人”这些标签，并清楚表明，自己将会为此付诸实践。
而这又把我们带向另一场反叛……
